The entire coefficient-two cubic image and the chosen lifts of all
corrected cup products can be removed simultaneously. Put
\[
 \mathcal O_2=\Z_2[\omega],\quad \omega^2+\omega+1=0,
 \quad \delta=1+2\omega,
 \quad \mathcal U=\{z\in1+4\mathcal O_2:z\bar z=1\}.
\]
Let \(\mathfrak t_2=T_{g,2}\) be the coefficient-two projection
of the corrected cubic CM homomorphism on \(M\).
Let \(\mathcal I\) consist of the odd-prime pairs \(\ell<q\),
and let \(\mathcal J\) consist of the odd primes other than three.

\begin{theorem}\label{thm:main-simultaneous-quadratic}
The integral logarithm defines an invariant character
\[
 \tau=\frac{\log\mathfrak t_2}{4\delta}\pmod2\in H^1(M,\F_2)^G,
 \qquad \tau|_K=\kappa_{2,3}+\psi_3.
\]
The joint map
\[
 (\mathfrak t_2,(\xi_{\ell,q})_{\mathcal I},(\nu_p)_{\mathcal J}):R_D
       \longrightarrow\mathcal U\times\prod_{\mathcal I}\F_2
                                   \times\prod_{\mathcal J}\F_2
\]
is surjective already on \(R_D\cap H\), and also on \(V_D\cap H\).
Define
\[
 \begin{split}
 R_{\rm quad}&=R_D\cap\ker\mathfrak t_2
          \cap\bigcap_{\mathcal I}\ker\xi_{\ell,q}
          \cap\bigcap_{\mathcal J}\ker\nu_p,\\
 V_{\rm quad}&=V_D\cap R_{\rm quad}.
 \end{split}
\]
These subgroups are closed and normal in \(G\), and
\[
 \begin{gathered}
 R_D/R_{\rm quad}\simeq V_D/V_{\rm quad}
   \simeq\mathcal U\times\prod_{\mathcal I}\F_2
                                   \times\prod_{\mathcal J}\F_2,\\
 G=R_{\rm quad}H,\qquad
 B\simeq R_{\rm quad}/
           \overline{[G,R_{\rm quad}](R_{\rm quad}\cap H)},\\
 B_{\max2}\simeq
 \left(V_{\rm quad}/
           \overline{[G,V_{\rm quad}](V_{\rm quad}\cap H)}\right)_{\max2}.
 \end{gathered}
\]
The two quotient descriptions also hold on arithmetic intersections.
Restriction from \(X_2\) to either residual subgroup has kernel
exactly \(\mathcal S=s(Y_2)\). Its image is precisely the space
of invariant quadratic characters on that subgroup vanishing on
its arithmetic intersection, canonically identified with
\(\operatorname{Hom}_{\rm cts}(B,\F_2)\).
\end{theorem}

This isolates the relative kernel after removing the complete known
arithmetic quadratic image. The subgroup \(V_{\rm quad}\) has
infinite index in \(V_D\), already arithmetically. The theorem
preserves the relative obstruction and does not compute its value.
The proofs are in Sections~\ref{sec:fr-cubic-coefficient-two},
\ref{sec:simultaneous-quadratic-normalization} and
\ref{sec:complete-quadratic-restriction}.
